\chapter{Objetivos}
\label{cap:objetivos}

A partir del problema planteado en el Capítulo~\ref{cap:introduccion} y del
marco revisado en el Capítulo~\ref{cap:marco-teorico}, este capítulo formula el
objetivo general del trabajo y lo desglosa en una serie de objetivos específicos
verificables que guían el desarrollo descrito en los capítulos posteriores.

%% =============================================================================
\section{Objetivo general}
\label{sec:objetivo-general}

El objetivo general del trabajo es diseñar e implementar un sistema
robótico de patrulla autónoma para estancias cerradas que cartografíe el
entorno, lo recorra de forma autónoma, detecte la presencia de personas como
incidencia y la notifique a una plataforma centralizada con la que el usuario
interactúe, todo ello articulado mediante una arquitectura orientada a eventos
que desacople la capa robótica del plano de control.

%% =============================================================================
\section{Objetivos específicos}
\label{sec:objetivos-especificos}

Para alcanzar el objetivo general se definen los siguientes objetivos
específicos:

\begin{enumerate}
  \item \textbf{OE1. Cartografía autónoma del entorno.} Dotar al robot de la
    capacidad de explorar de forma autónoma una estancia desconocida, construir
    su mapa mediante \gls{slam} y almacenarlo para su uso posterior.

  \item \textbf{OE2. Generación y ejecución de la ruta de patrulla.} Derivar del
    mapa una ruta de cobertura del espacio libre y recorrerla de manera cíclica
    con la pila de navegación \gls{nav2}, evitando obstáculos.

  \item \textbf{OE3. Orquestación del ciclo de patrulla.} Diseñar un orquestador,
    en forma de una máquina de estados finita (\acrshort{fsm}), que coordine las
    fases del ciclo operativo del robot (exploración y mapeado, planificación de
    la ruta, patrulla y retorno al punto de carga), de modo que el robot encadene
    de forma autónoma las capacidades anteriores.

  \item \textbf{OE4. Percepción y detección de incidencias en el entorno.}
    Integrar un módulo de visión por computador basado en \gls{yolo} que detecte
    personas durante la patrulla, grabe un clip del evento y lo registre como
    incidencia.

  \item \textbf{OE5. Arquitectura orientada a eventos.} Diseñar el sistema de
    comunicación basado en un bus de mensajería que desacople el robot del
    backend y canalice la telemetría, las incidencias y los comandos de control.

  \item \textbf{OE6. Backend y persistencia de datos.} Desarrollar un servicio
    de backend que consuma los eventos del robot, exponga una \gls{api}
    \gls{rest} y persista la información en los almacenes adecuados, incluyendo
    el vídeo de las incidencias y las series temporales de telemetría.

  \item \textbf{OE7. Seguridad y autenticación.} Implementar un mecanismo de
    autenticación tanto para el robot, mediante un esquema de confianza en el
    primer uso (\gls{tofu}) y tokens \gls{jwt}, como para los usuarios, con
    control de acceso basado en roles.

  \item \textbf{OE8. Aplicación de usuario.} Desarrollar una aplicación
    multiplataforma con \mbox{Flutter} desde la que el usuario controle la
    patrulla del robot y consulte el historial de incidencias y la telemetría.

  \item \textbf{OE9. Integración y validación funcional.} Integrar todos los
    subsistemas en una solución desplegable mediante contenedores y comprobar su
    funcionamiento conjunto.
\end{enumerate}

El grado de consecución de cada uno de estos objetivos se valora en el
Capítulo~\ref{cap:conclusiones}, a la luz de los resultados presentados en el
Capítulo~\ref{cap:resultados}.
